% %%%%%%%%%%%%%%%%%%%%%%%%%%%%%%%%%%%%%%%%%%%%%%%%%%%%%%%%%%%%%%%%%%%%%%%%%%%%%%%%%
%
%	cap1.tex 
%  -----------
%
%	Neste arquivo você escreve o capítulo. Deve sempre iniciar com o ambiente
%	\chapter{Nome do Capitulo}
%
% %%%%%%%%%%%%%%%%%%%%%%%%%%%%%%%%%%%%%%%%%%%%%%%%%%%%%%%%%%%%%%%%%%%%%%%%%%%%%%%%%
%
\chapter{Descrição do Trabalho e Escopo}
% ----------------------------------------------------------

\section{Objetivo Principal}

O desenvolvimento de uma ferramenta computacional para auxiliar engenheiros no
dimensionamento de vigas em concreto armado com aberturas na direção de sua largura.

\section{Objetivos Secundários}

O presente trabalho considera como secundários os seguintes objetivos:

\begin{enumerate}[label=\alph*)]
\item o desenvolvimento de um modelo estrutural baseado no método da rigidez para a obtenção dos esforços solicitantes;

\item desenvolver o roteiro de dimensionamento estrutural de vigas de concreto armado sob flexão, flexo-tração, flexo-compressão e corte;

\item implementar a solução em linguagem Python utilizando o paradigma de orientação a objetos;

\item desenvolver uma interface de usuário;

\item realizar verificações do modelo estrutural com softwares independentes;

\item comparação do dimensionamento em relação a métodos empíricos que constam na literatura;

\item realizar análises paramétricas referente a mudanças de posição ou de dimensão da abertura e seus impactos no dimensionamento estrutural.

\end{enumerate}

\section{Delimitações}

O presente trabalho terá as seguintes delimitações:

\begin{enumerate}[label=\alph*)]
\item apenas furos retangulares serão abordados;

\item apenas vigas retas sem efeito de torção;

\item análise estrutural através do modelo da rigidez;

\item serão consideradas apenas cargas distribuidas ao longo da estrutura;

\item vigas que respeitam a hipótese de seções planas (Euler-Bernoulli)

\end{enumerate}

\section{Delineamento}

Para atingir os resultados esperados, o trabalho será desenvolvido por meio das seguintes
etapas:

\begin{enumerate}[label=\alph*)]
\item revisão bibliográfica;

\item implementação do método da rigidez;

\item verificações parciais referente a análise estrutural;

\item implementação dos algoritmos para o dimensionamento das armaduras;

\item análises paramétricas;

\item verificação dos resultados;

\item considerações finais.

\end{enumerate}

A revisão bibliográfica consiste na pesquisa de normas, trabalhos acadêmicos e livros,
envolvendo o dimensionamento de vigas de concreto armado com aberturas, método da
rigidez e sua implementação computacional. O método da rigidez foi a ferramenta escolhida
para a obtenção dos esforços internos na estrutura, necessários para o dimensionamento. Para
verificar a consistência dos resultados, serão testados exemplos do livro Matrix Analysis of
Framed Structures de \citeonline{Weaver2018}, referência clássica na aplicação do método da
rigidez. O algoritmo para o dimensionamento das armaduras será planejado e elaborado
conforme os critérios e parâmetros da NBR 6118 \cite{ABNT_NBR6118_2026}. Na etapa de verificação dos
resultados, serão feitos testes finais no programa para avaliar seu funcionamento e
confiabilidade. Ao final serão feitas algumas análises paramétricas e considerações finais a
respeito do software. 
